\section{Problem Setting}\label{sec:setting}

We consider $N$ risky assets with daily simple returns $r_t\in\mathbb{R}^N$. At each
rebalancing date the manager observes the last $L$ return vectors, estimates a covariance
matrix $\hat\Sigma$, and chooses weights $w\in\mathbb{R}^N$. The mandate imposes a
polyhedral feasible set
\begin{equation}\label{eq:feasible}
\mathcal{W}(w^-)=\Bigl\{w:\ \mathbf{1}^\top w=1,\ \ \ell\le w\le u,\ \
\underline{c}_g\le \sum_{i\in g} w_i\le \bar c_g\ \ \forall g\in\mathcal{G},\ \
\lVert w-w^-\rVert_1\le\tau\Bigr\},
\end{equation}
where $\ell,u$ are per-asset bounds, $\mathcal{G}$ partitions the assets into groups
(sectors or asset classes) with caps $\bar c_g$ and floors $\underline{c}_g$, $w^-$ is the
drifted pre-trade portfolio and $\tau$ a turnover budget. All constraints are linear, so
$\mathcal{W}$ is a convex polytope. We assume it is non-empty. The turnover budget can in
principle conflict with the other constraints if price drift pushes the pre-trade
portfolio far outside the box or group limits; our implementation then relaxes the
turnover cap for that date only. In all experiments reported here this fallback was never
triggered: every CA-HRP portfolio satisfied all constraints, turnover included.

\section{Method}\label{sec:method}

\subsection{Hierarchical Risk Parity}
HRP \citep{lopezdeprado2016} proceeds in three steps. (i)~\emph{Tree clustering}: the
correlation matrix $\rho$ implied by $\hat\Sigma$ is mapped to the distance
$d_{ij}=\sqrt{(1-\rho_{ij})/2}$ and agglomerated with single linkage. (ii)~\emph{Quasi-
diagonalisation}: assets are ordered by the dendrogram leaves so that similar assets are
adjacent. (iii)~\emph{Recursive bisection}: starting from the full ordered list with weight
$W=1$, every cluster $C$ is split into a left and right half $C_L,C_R$; with
$V_S=\tilde w_S^\top\hat\Sigma_{SS}\tilde w_S$ the variance of the inverse-variance portfolio
$\tilde w_S\propto\operatorname{diag}(\hat\Sigma_{SS})^{-1}$ of a cluster $S$, the split
fraction is
\begin{equation}\label{eq:alpha}
\alpha_C = 1-\frac{V_{C_L}}{V_{C_L}+V_{C_R}},\qquad
W_{C_L}=\alpha_C W_C,\quad W_{C_R}=(1-\alpha_C)W_C .
\end{equation}
The recursion ends at single assets. No matrix is inverted and every split is a
two-asset risk-parity decision, which is the source of HRP's stability. Nothing in
\eqref{eq:alpha}, however, knows about $\mathcal{W}$.

\subsection{Capacity-aware recursive bisection (CA-HRP-B)}
Our first contribution makes the bisection itself respect per-asset bounds and, partially,
group caps, without any optimisation solver. For a set of assets $S$ define its
\emph{capacity interval}
\begin{equation}\label{eq:capacity}
\mathrm{lo}(S)=\sum_{i\in S}\ell_i,\qquad
\mathrm{hi}(S)=\sum_{g\in\mathcal{G}}\min\Bigl(\sum_{i\in S\cap g}u_i,\ \bar c_g\Bigr).
\end{equation}
$\mathrm{lo}(S)$ is the least weight that $S$ must hold and $\mathrm{hi}(S)$ an upper
bound on what it can hold. When a cluster $C$ carrying weight $W_C$ is split, the HRP
fraction \eqref{eq:alpha} is clipped to
\begin{equation}\label{eq:clip}
\alpha_C\in\Bigl[\tfrac{\max(\mathrm{lo}(C_L),\,W_C-\mathrm{hi}(C_R))}{W_C},\ \
\tfrac{\min(\mathrm{hi}(C_L),\,W_C-\mathrm{lo}(C_R))}{W_C}\Bigr].
\end{equation}

\begin{proposition}[Exact box feasibility]\label{prop:box}
If $\mathrm{lo}(\{1,\dots,N\})\le 1\le \mathrm{hi}(\{1,\dots,N\})$, the interval
\eqref{eq:clip} is non-empty at every split, and the resulting weights satisfy
$\mathbf{1}^\top w=1$ and $\ell_i\le w_i\le\min(u_i,\bar c_{g(i)})$ for all $i$.
\end{proposition}
\begin{proof}
We show by induction that every cluster satisfies the invariant
$\mathrm{lo}(C)\le W_C\le \mathrm{hi}(C)$. It holds at the root by assumption. Note that
$\mathrm{lo}$ is additive, $\mathrm{lo}(C)=\mathrm{lo}(C_L)+\mathrm{lo}(C_R)$, and
$\mathrm{hi}$ is sub-additive because $\min(a+b,c)\le\min(a,c)+\min(b,c)$ for
$a,b,c\ge0$, so $\mathrm{hi}(C)\le\mathrm{hi}(C_L)+\mathrm{hi}(C_R)$. The interval
\eqref{eq:clip} (multiplied by $W_C$) is non-empty iff the four inequalities
$\mathrm{lo}(C_L)\le\mathrm{hi}(C_L)$, $\mathrm{lo}(C_L)\le W_C-\mathrm{lo}(C_R)$,
$W_C-\mathrm{hi}(C_R)\le\mathrm{hi}(C_L)$ and $W_C-\mathrm{hi}(C_R)\le W_C-\mathrm{lo}(C_R)$
hold. The first and last follow from $\ell\le u$; the second from
$W_C\ge\mathrm{lo}(C)$ and additivity; the third from $W_C\le\mathrm{hi}(C)$ and
sub-additivity. Any $\alpha_C$ in the interval gives children with
$\mathrm{lo}(C_L)\le W_{C_L}\le\mathrm{hi}(C_L)$ and likewise for $C_R$, which preserves
the invariant. At a leaf $\{i\}$, $\mathrm{lo}=\ell_i$ and $\mathrm{hi}=\min(u_i,\bar
c_{g(i)})$, which proves the claim. Weights sum to one because each split conserves
$W_C$.
\end{proof}

CA-HRP-B costs $O(N^2)$ beyond HRP and needs no solver. It guarantees box bounds exactly,
but group constraints are only respected when a group lies inside one sub-tree, and group
floors and turnover are not handled. Those remaining constraints are enforced by a light
projection (next subsection, with the Euclidean metric replaced by a tree-weighted one;
details in the code). As our experiments show, the price of handling constraints
\emph{inside} the bisection is that excess weight cascades down a single branch of the
tree and can end up concentrated in volatile assets; this motivates the main method.

\subsection{Tracking-error projection (CA-HRP)}
The key observation is that HRP produces an excellent \emph{target} but has no notion of
distance between portfolios. Practitioners close the gap by clipping and renormalising,
which is not guaranteed to reach $\mathcal{W}$, or by a Euclidean projection
$\min_{w\in\mathcal{W}}\lVert w-w^{\mathrm{HRP}}\rVert_2^2$, which treats a unit of weight
moved from a short-term bond fund to an oil fund the same as one moved to another bond fund.
We instead measure distance in \emph{risk}: the ex-ante tracking-error variance of the
constrained portfolio relative to HRP \citep{roll1992},
\begin{equation}\label{eq:te}
w^{\mathrm{CA}}=\arg\min_{w\in\mathcal{W}(w^-)}\ (w-w^{\mathrm{HRP}})^\top\,\bar\Sigma\,
(w-w^{\mathrm{HRP}}),\qquad
\bar\Sigma=\frac{N\hat\Sigma_{\mathrm{LW}}}{\operatorname{tr}\hat\Sigma_{\mathrm{LW}}}+\epsilon I,
\end{equation}
where $\hat\Sigma_{\mathrm{LW}}$ is the Ledoit--Wolf shrinkage estimator
\citep{ledoit2004}, $w^{\mathrm{HRP}}$ is HRP computed on $\hat\Sigma_{\mathrm{LW}}$, and
the small ridge $\epsilon=10^{-4}$ guarantees strict convexity.

\begin{proposition}[Properties of the projection]\label{prop:te}
(a) If $\mathcal{W}(w^-)\neq\emptyset$, problem \eqref{eq:te} has a unique solution.
(b) If $w^{\mathrm{HRP}}\in\mathcal{W}(w^-)$ then $w^{\mathrm{CA}}=w^{\mathrm{HRP}}$:
CA-HRP never alters HRP when HRP already complies.
(c) Among all feasible portfolios, $w^{\mathrm{CA}}$ minimises the ex-ante variance of
the return difference to HRP, so it is the feasible portfolio whose returns are most
correlated with, and least volatile around, the HRP portfolio.
(d) Let $A w\le b$ collect the active inequality constraints and $\mathbf 1^\top w=1$ the
budget. The KKT conditions give
$w^{\mathrm{CA}}=w^{\mathrm{HRP}}-\tfrac12\bar\Sigma^{-1}(A^\top\lambda+\nu\mathbf{1})$
with $\lambda\ge0$: the correction is a \emph{hedge}. Weight removed from a capped
asset is redistributed in proportion to the rows of $\bar\Sigma^{-1}$, i.e.\ towards the
assets that best replicate the capped asset's risk.
\end{proposition}
\begin{proof}
(a) $\bar\Sigma\succ0$, so the objective is strictly convex and $\mathcal{W}$ is a non-empty
closed convex set; a unique minimiser exists. (b) The objective is non-negative and zero at
$w^{\mathrm{HRP}}$. (c) $(w-w^{\mathrm{HRP}})^\top\hat\Sigma(w-w^{\mathrm{HRP}})$ is the
variance of $(w-w^{\mathrm{HRP}})^\top r$ under the estimated covariance; the trace
normalisation and ridge only rescale and regularise it. (d) Stationarity of the
Lagrangian $(w-w^{\mathrm{HRP}})^\top\bar\Sigma(w-w^{\mathrm{HRP}})+\lambda^\top(Aw-b)
+\nu(\mathbf{1}^\top w-1)$, with the $\ell_1$ turnover constraint written in its standard
linear epigraph form.
\end{proof}

Part (d) explains the behaviour we observe in the multi-asset universe: when HRP wants far
more than the 40\% asset-class cap in bonds, the Euclidean projection spreads the excess
almost uniformly (about 5\% to each non-bond fund), while the tracking-error projection
re-shapes the portfolio as a hedge: inside the bond budget it moves from short-term to
intermediate and long Treasuries (IEF, TLT), whose negative correlation with equities
offsets the risk of the assets that must now be held, and outside the bond budget it
favours gold and broad commodities (Section~\ref{sec:results}). The constrained portfolio
keeps the defensive character of HRP, with markedly lower volatility than the Euclidean
alternative.

\begin{algorithm}[t]
\caption{Constraint-Aware HRP (one rebalancing date)}\label{alg:cahrp}
\begin{algorithmic}[1]
\Require returns window $R\in\mathbb{R}^{L\times N}$, feasible set $\mathcal{W}$, drifted weights $w^-$
\State $\hat\Sigma_{\mathrm{LW}}\gets$ Ledoit--Wolf shrinkage covariance of $R$
\State $\rho\gets$ correlation of $\hat\Sigma_{\mathrm{LW}}$;\ \ $d_{ij}\gets\sqrt{(1-\rho_{ij})/2}$
\State $\pi\gets$ leaf order of the single-linkage dendrogram of $d$
\State $w^{\mathrm{HRP}}\gets$ recursive bisection \eqref{eq:alpha} on $\pi$
\Comment{optionally clipped by \eqref{eq:clip}}
\State $w^{\mathrm{CA}}\gets$ solve the convex QP \eqref{eq:te} over $\mathcal{W}(w^-)$
\State \Return $w^{\mathrm{CA}}$
\end{algorithmic}
\end{algorithm}

\paragraph{Complexity.} Steps 1--4 cost $O(LN^2+N^2\log N)$. Step~5 is a QP with $N$
variables (plus $N$ auxiliary variables for the turnover epigraph) and $N+|\mathcal{G}|$
linear constraints; with an interior-point solver \citep{diamond2016} it takes a few
milliseconds for $N\le 50$. A full 20-year monthly backtest of CA-HRP runs in under
six seconds on a laptop; Section~\ref{sec:extended} reports timings up to $N=400$.

\subsection{Two further properties: noise damping and the price of compliance}
\label{sec:props}

\begin{proposition}[Noise damping]\label{prop:nonexp}
Fix the metric $\bar\Sigma\succ0$ and the feasible set $\mathcal{W}$, and let
$P(t)=\arg\min_{w\in\mathcal{W}}\lVert w-t\rVert_{\bar\Sigma}$ with
$\lVert x\rVert_{\bar\Sigma}^2=x^\top\bar\Sigma x$. Then for any two targets $t_1,t_2$,
$\lVert P(t_1)-P(t_2)\rVert_{\bar\Sigma}\le\lVert t_1-t_2\rVert_{\bar\Sigma}$.
\end{proposition}
\begin{proof}
$\langle x,y\rangle_{\bar\Sigma}=x^\top\bar\Sigma y$ is an inner product, so $P$ is the
metric projection onto a closed convex set in a Hilbert space, which is firmly
non-expansive \citep{bauschke2017}: $\lVert P t_1-Pt_2\rVert^2\le\langle Pt_1-Pt_2,\,t_1-t_2\rangle\le
\lVert Pt_1-Pt_2\rVert\,\lVert t_1-t_2\rVert$ by the variational inequality characterising
projections and the Cauchy--Schwarz inequality.
\end{proof}

Proposition~\ref{prop:nonexp} means that whatever estimation noise the HRP target carries,
the projection can only shrink it, measured in the same risk metric. The clip-and-renormalise
heuristic has no such property; Section~\ref{sec:round3} shows that in practice it often
\emph{amplifies} noise.

\begin{proposition}[Exact price-of-compliance identity]\label{prop:kkt}
Let $f^\star=(w^{\mathrm{CA}}-t)^\top\bar\Sigma(w^{\mathrm{CA}}-t)$ be the optimal value of
\eqref{eq:te} with linear constraints $a_k^\top w\le b_k$ and multipliers
$\lambda_k\ge0$. Then
\begin{equation}\label{eq:kkt}
f^\star=\tfrac12\sum_k\lambda_k\,(a_k^\top t-b_k),
\end{equation}
so the cost of compliance splits exactly over constraints, each term being the shadow price
of a constraint times the amount by which the HRP target $t$ violates it.
\end{proposition}
\begin{proof}
Stationarity gives $2\bar\Sigma d=-(\sum_k\lambda_k a_k+\nu\mathbf1)$ with $d=w^{\mathrm{CA}}-t$.
Multiply by $d^\top/2$: $f^\star=-\tfrac12\sum_k\lambda_k a_k^\top d$, because
$\mathbf1^\top d=0$. Complementary slackness gives
\[
\lambda_k a_k^\top w^{\mathrm{CA}}=\lambda_k b_k,\qquad\text{hence}\qquad
-\lambda_k a_k^\top d=\lambda_k(a_k^\top t-b_k).
\]
\end{proof}

Individual terms of \eqref{eq:kkt} can be negative: a cap that the target already satisfies
but that binds at the optimum contributes negatively. To report the cost by \emph{family} of
constraints (asset caps; group caps and floors; turnover) we therefore use the Shapley value \citep{shapley1953}
of the cooperative game $v(S)=f^\star_S$, the optimal value when only the families in $S$
are imposed. It is order-independent and adds up exactly to $f^\star$, and it costs $2^3=8$
small solves per date. The $\ell_1$ turnover cap is not linear in $w$, but in its standard
linear epigraph form the identity still holds; we compute the turnover term as
$f^\star$ minus the linear terms and use \eqref{eq:kkt} as an exactness check on every date
(the remainder must vanish whenever the turnover cap is slack).

\subsection{Regime-robust projection (CA-HRP-R)}\label{sec:robust}
The tracking-error metric in \eqref{eq:te} trusts a single covariance estimate. When
correlations shift, the hedge chosen by the projection can fail. We therefore also study a
\emph{regime-robust} projection that tracks HRP under the worst of $K$ covariance
scenarios $\{\Sigma_k\}_{k=1}^K$:
\begin{equation}\label{eq:robust}
w^{\mathrm{R}}=\arg\min_{w\in\mathcal{W}(w^-)}\ \max_{k=1,\dots,K}\
(w-w^{\mathrm{HRP}})^\top\bar\Sigma_k(w-w^{\mathrm{HRP}}),
\end{equation}
where every $\bar\Sigma_k$ uses the same scale factor as the baseline estimate, so that
relative risk levels across regimes are preserved. We consider three scenarios: the
1-year Ledoit--Wolf estimate, a 3-month (63-day) Ledoit--Wolf estimate capturing the
current regime, and a ``hedges-fail'' stress covariance
$D\lvert\rho\rvert D$ in which every correlation is replaced by its absolute value.

\begin{proposition}[Robust projection]\label{prop:robust}
Problem \eqref{eq:robust} is convex, has a unique solution, and returns $w^{\mathrm{HRP}}$
whenever it is feasible. Its optimal value $t^\star$ bounds the ex-ante tracking-error
variance under \emph{every} covariance matrix in the convex hull
$\operatorname{conv}\{\bar\Sigma_1,\dots,\bar\Sigma_K\}$:
$(w^{\mathrm{R}}-w^{\mathrm{HRP}})^\top\Sigma(w^{\mathrm{R}}-w^{\mathrm{HRP}})\le t^\star$.
\end{proposition}
\begin{proof}
Each term is a strictly convex quadratic, and the pointwise maximum of strictly convex
functions is strictly convex; with a non-empty convex feasible set this gives existence
and uniqueness, and the objective is zero only at $w^{\mathrm{HRP}}$. For
$\Sigma=\sum_k\lambda_k\bar\Sigma_k$ with $\lambda$ in the simplex and $d=w-w^{\mathrm{HRP}}$,
$d^\top\Sigma d=\sum_k\lambda_k d^\top\bar\Sigma_k d\le\max_k d^\top\bar\Sigma_k d$, because
the tracking-error variance is linear in $\Sigma$.
\end{proof}

Problem \eqref{eq:robust} is a second-order cone program (epigraph form) and solves as fast
as \eqref{eq:te}. The choice of scenarios was made on the validation period only
(Appendix~\ref{app:design}): the variant using the 1-year and 3-month regimes was selected.
The hedges-fail stress scenario was rejected at that stage because it lowered validation
performance markedly, even though the held-out period later contained exactly such a
regime (2022). We report the frozen choice regardless.
